

\documentclass{beamer}

\hypersetup{pdfpagemode=FullScreen}
\setbeamertemplate{navigation symbols}{}

\usepackage[english]{babel}
\usepackage[utf8]{inputenc}
\usepackage[T1]{fontenc}

\usepackage{amsmath,amssymb}
\usepackage{graphicx}
\usepackage{booktabs}
\usepackage{array}
\usepackage{xcolor}
\usepackage{hyperref}
\usepackage{mathptmx}

\mode<presentation>{
\usetheme{Madrid}
\usecolortheme{whale}
}

\setbeamercovered{transparent}
\setbeamerfont{title}{size=\large}
\setbeamerfont{date}{size=\tiny}
\setbeamertemplate{itemize items}[ball]
\setbeamertemplate{itemize subitem}[circle]
\setbeamertemplate{itemize subsubitem}[triangle]
\setbeamertemplate{caption}[numbered]

%%%%%%%%%%%%%%%%%%%%%%%%%%%%%%%%%%%%%%%%%%%%%%%%%%%%%%%%%%%%%%%%%%%%%%%%%%
% TITLE PAGE
%%%%%%%%%%%%%%%%%%%%%%%%%%%%%%%%%%%%%%%%%%%%%%%%%%%%%%%%%%%%%%%%%%%%%%%%%%

\title[History of AI \& Deep Learning]
{History of Artificial Intelligence \& Deep Learning}


%%%%%%%%%%%%%%%%%%%%%%%%%%%%%%%%%%%%%%%%%%%%%%%%%%%%%%%%%%%%%%%%%%%%%%%%%%

\begin{document}

\begin{frame}
\titlepage
\end{frame}

\begin{frame}{Contents}
\tableofcontents
\end{frame}
\author[Vatsav Reddy]
{
{\Large Vatsav Reddy}\\
\vspace{0.2cm}
{\small Under the Guidance of}\\
{\small Prof.\ Praveen Parchuri}
}

\institute[IIITH]
{
{\normalsize International Institute of Information Technology, Hyderabad}
}
\AtBeginSection[]{
\begin{frame}{Current Section}
\tableofcontents[currentsection]
\end{frame}
}

%%%%%%%%%%%%%%%%%%%%%%%%%%%%%%%%%%%%%%%%%%%%%%%%%%%%%%%%%%%%%%%%%%%%%%%%%%
\section{Lecture 1-1: Biological Neurons}
%%%%%%%%%%%%%%%%%%%%%%%%%%%%%%%%%%%%%%%%%%%%%%%%%%%%%%%%%%%%%%%%%%%%%%%%%%

\begin{frame}{Course Overview}

\begin{block}{Topics Covered}
\begin{itemize}
\item Deep Generative Models
\item Restricted Boltzmann Machines
\item Variational Autoencoders (VAEs)
\item Autoregressive Models
\item Generative Adversarial Networks (GANs)
\end{itemize}
\end{block}

\vspace{0.2cm}

\begin{block}{Main Architectures}
CNNs, RNNs, LSTMs, Attention Mechanisms
\end{block}

\vspace{0.2cm}

\begin{exampleblock}{Reference}
J.\ Schmidhuber, \textit{Deep Learning: An Overview}
\end{exampleblock}

\end{frame}

%%%%%%%%%%%%%%%%%%%%%%%%%%%%%%%%%%%%%%%%%%%%%%%%%%%%%%%%%%%%%%%%%%%%%%%%%%

\begin{frame}{Reticular Theory vs.\ Neuron Doctrine}

\begin{columns}

\column{0.48\textwidth}

\begin{block}{Reticular Theory (1871--1873)}
\begin{itemize}
\item Proposed by Joseph von Gerlach
\item Nervous system viewed as a single continuous network
\item Camillo Golgi developed staining techniques
\item Golgi supported Reticular Theory
\end{itemize}
\end{block}

\column{0.48\textwidth}

\begin{alertblock}{Neuron Doctrine (1888--1891)}
\begin{itemize}
\item Proposed by Santiago Ramón y Cajal
\item Nervous system made of discrete neurons
\item Neurons communicate through synapses
\item Term ``neuron'' coined by Waldeyer-Hartz
\end{itemize}
\end{alertblock}

\end{columns}

\vspace{0.2cm}

\begin{exampleblock}{1906 Nobel Prize}
Golgi and Cajal jointly received the Nobel Prize in Physiology/Medicine despite supporting opposing theories.
\end{exampleblock}

\end{frame}

%%%%%%%%%%%%%%%%%%%%%%%%%%%%%%%%%%%%%%%%%%%%%%%%%%%%%%%%%%%%%%%%%%%%%%%%%%

\begin{frame}{The Final Word}

\begin{block}{Electron Microscopy (1950s)}
Electron microscopy conclusively confirmed the Neuron Doctrine by showing that neurons are distinct individual cells connected through synapses.
\end{block}

\end{frame}

%%%%%%%%%%%%%%%%%%%%%%%%%%%%%%%%%%%%%%%%%%%%%%%%%%%%%%%%%%%%%%%%%%%%%%%%%%
\section{Lecture 1-2: From Spring to Winter of AI}
%%%%%%%%%%%%%%%%%%%%%%%%%%%%%%%%%%%%%%%%%%%%%%%%%%%%%%%%%%%%%%%%%%%%%%%%%%

\begin{frame}{McCulloch-Pitts Neuron (1943)}

\begin{itemize}
\item Proposed by Warren McCulloch and Walter Pitts
\item Simplified mathematical model of a biological neuron
\item Inputs: $x_1,x_2,\ldots,x_n$
\item Output: binary value $y \in \{0,1\}$
\end{itemize}

\vspace{0.3cm}

\begin{block}{Threshold Logic}

\[
y =
f\left(
\sum_{i=1}^{n} w_i x_i
\right)
=
\begin{cases}
1 & \text{if } \sum w_i x_i \ge \theta \\
0 & \text{otherwise}
\end{cases}
\]

\end{block}

\end{frame}

%%%%%%%%%%%%%%%%%%%%%%%%%%%%%%%%%%%%%%%%%%%%%%%%%%%%%%%%%%%%%%%%%%%%%%%%%%

\begin{frame}{The Perceptron (1957--1958)}

\begin{itemize}
\item ``Artificial Intelligence'' term coined in 1956
\item Frank Rosenblatt proposed the Perceptron
\item Inputs connected using learnable weights
\item Could learn simple decision boundaries
\item Received enormous media attention
\end{itemize}

\vspace{0.2cm}

\begin{exampleblock}{Historical Quote}
``The perceptron may eventually be able to learn, make decisions, and translate languages.''
\end{exampleblock}

\vspace{0.2cm}

\begin{block}{First Multilayer Perceptrons}
Ivakhnenko et al.\ (1965--1968) proposed early multilayer neural networks.
\end{block}

\end{frame}

%%%%%%%%%%%%%%%%%%%%%%%%%%%%%%%%%%%%%%%%%%%%%%%%%%%%%%%%%%%%%%%%%%%%%%%%%%

\begin{frame}{Perceptron Limitations and AI Winter}

\begin{alertblock}{Minsky \& Papert (1968)}
\begin{itemize}
\item Perceptrons cannot solve XOR
\item XOR is not linearly separable
\item Funding for connectionist AI decreased
\item Led to the first AI Winter
\end{itemize}
\end{alertblock}

\vspace{0.2cm}

\begin{block}{Backpropagation (1986)}
\begin{itemize}
\item Rediscovered several times during 1960s--1970s
\item Paul Werbos used it for neural networks in 1982
\item Popularized by Rumelhart, Hinton, and Williams (1986)
\end{itemize}
\end{block}

\end{frame}

%%%%%%%%%%%%%%%%%%%%%%%%%%%%%%%%%%%%%%%%%%%%%%%%%%%%%%%%%%%%%%%%%%%%%%%%%%
\section{Lecture 1-3: The Deep Revival}
%%%%%%%%%%%%%%%%%%%%%%%%%%%%%%%%%%%%%%%%%%%%%%%%%%%%%%%%%%%%%%%%%%%%%%%%%%

\begin{frame}{From AI Winter to Deep Revival}

\begin{itemize}
\item 1847: Cauchy introduced Gradient Descent
\item 1989: Universal Approximation Theorem
\item 1989--2006: Second AI Winter
\end{itemize}

\vspace{0.2cm}

\begin{block}{Universal Approximation Theorem}
A feedforward neural network with a single hidden layer can approximate any continuous function to arbitrary precision.
\end{block}

\vspace{0.2cm}

\begin{exampleblock}{Deep Revival (2006)}
Hinton and Salakhutdinov introduced effective pretraining methods for deep autoencoders.
\end{exampleblock}

\end{frame}

%%%%%%%%%%%%%%%%%%%%%%%%%%%%%%%%%%%%%%%%%%%%%%%%%%%%%%%%%%%%%%%%%%%%%%%%%%

\begin{frame}{Unsupervised Pretraining and GPU Era}

\begin{itemize}

\item 1991--1993: Schmidhuber explored unsupervised pretraining

\item 2007--2009: Greedy layer-wise pretraining became popular

\item 2009: Graves et al.\ achieved Arabic handwriting recognition

\item 2010: Dahl et al.\ reduced speech recognition error:
\[
23.2\% \rightarrow 16.0\%
\]

\item GPUs dramatically accelerated backpropagation

\item 2011: Traffic sign recognition achieved 0.56\% error rate

\end{itemize}

\end{frame}

%%%%%%%%%%%%%%%%%%%%%%%%%%%%%%%%%%%%%%%%%%%%%%%%%%%%%%%%%%%%%%%%%%%%%%%%%%

\begin{frame}{ImageNet Challenge}

\begin{center}

\begin{tabular}{lcc}
\toprule
\textbf{Network} & \textbf{Top-5 Error} & \textbf{Layers} \\
\midrule
AlexNet & 16.4\% & 8 \\
ZFNet & 11.2\% & 8 \\
VGGNet & 7.3\% & 19 \\
GoogLeNet & 6.7\% & 22 \\
ResNet (2015) & 3.6\% & 152 \\
\bottomrule
\end{tabular}

\end{center}

\end{frame}

%%%%%%%%%%%%%%%%%%%%%%%%%%%%%%%%%%%%%%%%%%%%%%%%%%%%%%%%%%%%%%%%%%%%%%%%%%
\section{Lecture 1-4: From Cats to CNNs}
%%%%%%%%%%%%%%%%%%%%%%%%%%%%%%%%%%%%%%%%%%%%%%%%%%%%%%%%%%%%%%%%%%%%%%%%%%

\begin{frame}{Hubel and Wiesel Experiment (1959)}

\begin{block}{Key Discovery}
Each neuron responds only to stimuli in a specific receptive field of visual space.
\end{block}

\vspace{0.2cm}

\begin{itemize}
\item ``Neurons that fire together wire together''
\item Inspired hierarchical visual processing models
\end{itemize}

\end{frame}

%%%%%%%%%%%%%%%%%%%%%%%%%%%%%%%%%%%%%%%%%%%%%%%%%%%%%%%%%%%%%%%%%%%%%%%%%%

\begin{frame}{Neocognitron and CNNs}

\begin{block}{Neocognitron (1980)}
Kunihiko Fukushima proposed the Neocognitron for handwritten character recognition.
\end{block}

\vspace{0.2cm}

\begin{itemize}
\item Late 1980s--1990s: CNNs popularized by Yann LeCun
\item LeNet-5 (1998) achieved success on MNIST
\item CNNs later enabled large-scale image recognition
\end{itemize}

\vspace{0.2cm}

\begin{exampleblock}{Max Pooling}
\[
J(x,y)=\max_{m,n} I(x+m,y+n)
\]
\end{exampleblock}

\end{frame}

%%%%%%%%%%%%%%%%%%%%%%%%%%%%%%%%%%%%%%%%%%%%%%%%%%%%%%%%%%%%%%%%%%%%%%%%%%
\section{Lecture 1-5: Faster, Higher, Stronger}
%%%%%%%%%%%%%%%%%%%%%%%%%%%%%%%%%%%%%%%%%%%%%%%%%%%%%%%%%%%%%%%%%%%%%%%%%%

\begin{frame}{Better Optimization Techniques}

\begin{block}{Optimization Advances (2012--2016)}
\begin{itemize}
\item Faster convergence
\item Better generalization
\item More stable training
\item Nesterov Momentum improved SGD
\end{itemize}
\end{block}

\vspace{0.2cm}

\begin{center}

\begin{tabular}{ll}
\toprule
\textbf{Year} & \textbf{Optimizer} \\
\midrule
2001 & AdaGrad \\
2012 & RMSProp \\
2014 & Adam \\
2015 & Eve \\
2018 & Beyond Adam \\
\bottomrule
\end{tabular}

\end{center}

\end{frame}

%%%%%%%%%%%%%%%%%%%%%%%%%%%%%%%%%%%%%%%%%%%%%%%%%%%%%%%%%%%%%%%%%%%%%%%%%%
\section{Lecture 1-6: The Curious Case of Sequences}
%%%%%%%%%%%%%%%%%%%%%%%%%%%%%%%%%%%%%%%%%%%%%%%%%%%%%%%%%%%%%%%%%%%%%%%%%%

\begin{frame}{Why Sequences Matter}

\begin{block}{Sequences Are Everywhere}
\begin{itemize}
\item Time series
\item Speech
\item Music
\item Text
\item Video
\end{itemize}
\end{block}

\vspace{0.2cm}

\begin{exampleblock}{Hopfield Networks (1982)}
Associative memory systems capable of recalling patterns from noisy inputs.
\end{exampleblock}

\end{frame}

%%%%%%%%%%%%%%%%%%%%%%%%%%%%%%%%%%%%%%%%%%%%%%%%%%%%%%%%%%%%%%%%%%%%%%%%%%

\begin{frame}{Jordan and Elman Networks}

\begin{columns}

\column{0.48\textwidth}

\begin{block}{Jordan Network (1986)}
Output state is fed back into future time steps.
\end{block}

\column{0.48\textwidth}

\begin{block}{Elman Network (1990)}
Hidden state is propagated across time steps.
\end{block}

\end{columns}

\end{frame}

%%%%%%%%%%%%%%%%%%%%%%%%%%%%%%%%%%%%%%%%%%%%%%%%%%%%%%%%%%%%%%%%%%%%%%%%%%

\begin{frame}{RNN Problems and LSTMs}

\begin{alertblock}{Exploding and Vanishing Gradients}
Hochreiter and Bengio demonstrated the difficulty of training RNNs over long sequences.
\end{alertblock}

\vspace{0.2cm}

\begin{block}{Long Short-Term Memory (LSTM) (1997)}
LSTMs solved long-term dependency problems using gated memory cells.
\end{block}

\vspace{0.2cm}

\begin{exampleblock}{Sequence-to-Sequence Learning}
\begin{itemize}
\item Sutskever et al.\ (2014)
\item Bahdanau et al.\ (2015) introduced Attention
\end{itemize}
\end{exampleblock}

\end{frame}

%%%%%%%%%%%%%%%%%%%%%%%%%%%%%%%%%%%%%%%%%%%%%%%%%%%%%%%%%%%%%%%%%%%%%%%%%%

\begin{frame}{Attention and Reinforcement Learning}

\begin{block}{Schmidhuber \& Huber (1991)}
Used reinforcement learning to determine where a model should focus attention.
\end{block}

\end{frame}

%%%%%%%%%%%%%%%%%%%%%%%%%%%%%%%%%%%%%%%%%%%%%%%%%%%%%%%%%%%%%%%%%%%%%%%%%%
\section{Lecture 1-7: Beating Humans at Their Own Games}
%%%%%%%%%%%%%%%%%%%%%%%%%%%%%%%%%%%%%%%%%%%%%%%%%%%%%%%%%%%%%%%%%%%%%%%%%%

\begin{frame}{AI Beats Humans at Games}

\begin{itemize}

\item 2015: Deep Q-Networks (DQN) mastered Atari games

\item 2016: AlphaGo defeated world champion Lee Sedol

\item 2016: DeepStack defeated professional poker players

\item 2017: OpenAI Five played Defense of the Ancients 2 (Dota 2)

\end{itemize}

\end{frame}

%%%%%%%%%%%%%%%%%%%%%%%%%%%%%%%%%%%%%%%%%%%%%%%%%%%%%%%%%%%%%%%%%%%%%%%%%%
\section{Lecture 1-8: The Modern Era}
%%%%%%%%%%%%%%%%%%%%%%%%%%%%%%%%%%%%%%%%%%%%%%%%%%%%%%%%%%%%%%%%%%%%%%%%%%

\begin{frame}{Language Modeling and Translation}

\begin{columns}

\column{0.48\textwidth}

\begin{block}{Language Modeling}
\begin{itemize}
\item Mikolov et al.\ 2010
\item Kiros et al.\ 2015
\item Kim et al.\ 2015
\end{itemize}
\end{block}

\begin{block}{Machine Translation}
\begin{itemize}
\item Kalchbrenner et al.\ 2013
\item Cho et al.\ 2014
\item Bahdanau et al.\ 2015
\item Luong et al.\ 2015
\item Sutskever et al.\ 2014
\end{itemize}
\end{block}

\column{0.48\textwidth}

\begin{block}{Speech Recognition}
\begin{itemize}
\item Hinton et al.\ 2012
\item Graves et al.\ 2013
\item Chorowski et al.\ 2015
\item Sak et al.\ 2015
\end{itemize}
\end{block}

\end{columns}

\end{frame}

%%%%%%%%%%%%%%%%%%%%%%%%%%%%%%%%%%%%%%%%%%%%%%%%%%%%%%%%%%%%%%%%%%%%%%%%%%

\begin{frame}{Conversation and Question Answering}

\begin{columns}

\column{0.48\textwidth}

\begin{block}{Conversation Modeling}
\begin{itemize}
\item Shang et al.\ 2015
\item Vinyals et al.\ 2015
\item Lowe et al.\ 2015
\item Serban et al.\ 2016
\item Bordes et al.\ 2017
\end{itemize}
\end{block}

\column{0.48\textwidth}

\begin{block}{Question Answering}
\begin{itemize}
\item Hermann et al.\ 2015
\item Chen et al.\ 2016
\item Xiong et al.\ 2016
\item Dhingra et al.\ 2017
\item Wang et al.\ 2017
\end{itemize}
\end{block}

\end{columns}

\end{frame}

%%%%%%%%%%%%%%%%%%%%%%%%%%%%%%%%%%%%%%%%%%%%%%%%%%%%%%%%%%%%%%%%%%%%%%%%%%

\begin{frame}{Object Detection and Recognition}

\begin{itemize}

\item Semantic Segmentation (Long et al.\ 2015)

\item Faster R-CNN (Ren et al.\ 2015)

\item Inside-Outside Net (Bell et al.\ 2015)

\item YOLO9000 (Redmon et al.\ 2016)

\item R-FCN (Dai et al.\ 2016)

\item Mask R-CNN (2017)

\item Video Object Segmentation (Caelles et al.\ 2017)

\end{itemize}

\end{frame}

%%%%%%%%%%%%%%%%%%%%%%%%%%%%%%%%%%%%%%%%%%%%%%%%%%%%%%%%%%%%%%%%%%%%%%%%%%

\begin{frame}{Visual Question Answering}

\begin{itemize}
\item Agrawal et al.\ 2015
\item Yu et al.\ 2017
\item Johnson et al.\ 2017
\item Ben-Younes et al.\ 2017
\item Malinowski et al.\ 2017
\item Kazemi et al.\ 2017
\end{itemize}

\end{frame}

%%%%%%%%%%%%%%%%%%%%%%%%%%%%%%%%%%%%%%%%%%%%%%%%%%%%%%%%%%%%%%%%%%%%%%%%%%

\begin{frame}{Video Question Answering}

\begin{itemize}
\item Tapaswi et al.\ 2016
\item Zeng et al.\ 2016
\item Zhao et al.\ 2017
\item Jang et al.\ 2017
\item Xu Hongyang et al.\ 2017
\item Maharaj et al.\ 2017
\end{itemize}

\end{frame}

%%%%%%%%%%%%%%%%%%%%%%%%%%%%%%%%%%%%%%%%%%%%%%%%%%%%%%%%%%%%%%%%%%%%%%%%%%

\begin{frame}{Image and Video Captioning}

\begin{columns}

\column{0.48\textwidth}

\begin{block}{Image Captioning}
\begin{itemize}
\item Mao et al.\ 2014
\item Kiros et al.\ 2015
\item Donahue et al.\ 2015
\item Vinyals et al.\ 2015
\item Karpathy et al.\ 2015
\end{itemize}
\end{block}

\column{0.48\textwidth}

\begin{block}{Video Captioning}
\begin{itemize}
\item Donahue et al.\ 2014
\item Venugopalan et al.\ 2014
\item Pan et al.\ 2015
\item Yao et al.\ 2015
\item Rohrbach et al.\ 2015
\end{itemize}
\end{block}

\end{columns}

\end{frame}

%%%%%%%%%%%%%%%%%%%%%%%%%%%%%%%%%%%%%%%%%%%%%%%%%%%%%%%%%%%%%%%%%%%%%%%%%%

\begin{frame}{Generative Models}

\begin{block}{Generating Realistic Photos}
\begin{itemize}
\item Variational Autoencoders --- Kingma et al.\ 2013
\item GANs --- Goodfellow et al.\ 2014
\item Plug and Play Generative Networks --- Nguyen et al.\ 2016
\item Progressive Growing of GANs --- Karras et al.\ 2017
\end{itemize}
\end{block}

\vspace{0.2cm}

\begin{block}{Generating Audio}
WaveNet --- van den Oord et al.\ 2016
\end{block}

\end{frame}

%%%%%%%%%%%%%%%%%%%%%%%%%%%%%%%%%%%%%%%%%%%%%%%%%%%%%%%%%%%%%%%%%%%%%%%%%%

\begin{frame}{Transformers (2017)}

\begin{block}{Attention Is All You Need}
Vaswani et al.\ introduced the Transformer architecture based entirely on self-attention mechanisms.
\end{block}

\vspace{0.2cm}

\begin{itemize}
\item Removed recurrence
\item Highly parallelizable
\item Became foundation of modern LLMs
\end{itemize}

\end{frame}

%%%%%%%%%%%%%%%%%%%%%%%%%%%%%%%%%%%%%%%%%%%%%%%%%%%%%%%%%%%%%%%%%%%%%%%%%%
\section{Lecture 1-9: Need for Sanity}
%%%%%%%%%%%%%%%%%%%%%%%%%%%%%%%%%%%%%%%%%%%%%%%%%%%%%%%%%%%%%%%%%%%%%%%%%%

\begin{frame}{Why Does Deep Learning Work So Well?}

\begin{alertblock}{Open Problems}
\begin{itemize}
\item High capacity and overfitting
\item Numerical instability
\item Exploding and vanishing gradients
\item Sharp minima
\item Adversarial examples
\end{itemize}
\end{alertblock}

\vspace{0.2cm}

\begin{block}{Current Direction}
\begin{itemize}
\item Explainability
\item Robustness
\item Theoretical understanding
\item Reliable AI systems
\end{itemize}
\end{block}

\end{frame}

%%%%%%%%%%%%%%%%%%%%%%%%%%%%%%%%%%%%%%%%%%%%%%%%%%%%%%%%%%%%%%%%%%%%%%%%%%
\section{Summary}
%%%%%%%%%%%%%%%%%%%%%%%%%%%%%%%%%%%%%%%%%%%%%%%%%%%%%%%%%%%%%%%%%%%%%%%%%%

\begin{frame}{Timeline Summary}

\begin{itemize}

\item 1871 -- Reticular Theory

\item 1888 -- Neuron Doctrine

\item 1943 -- McCulloch-Pitts Neuron

\item 1957 -- Perceptron

\item 1968 -- AI Winter

\item 1980 -- Neocognitron

\item 1986 -- Backpropagation

\item 1989 -- Universal Approximation Theorem

\item 1997 -- LSTM

\item 2006 -- Deep Revival

\item 2012 -- AlexNet and GPU Revolution

\item 2014 -- Seq2Seq and GANs

\item 2015 -- DQN and ResNet

\item 2016 -- AlphaGo

\item 2017 -- Transformers

\end{itemize}

\end{frame}

%%%%%%%%%%%%%%%%%%%%%%%%%%%%%%%%%%%%%%%%%%%%%%%%%%%%%%%%%%%%%%%%%%%%%%%%%%

\begin{frame}
\Huge{\centerline{Thank You}}
\vspace{0.5cm}
\Large{\centerline{Questions?}}
\end{frame}

%%%%%%%%%%%%%%%%%%%%%%%%%%%%%%%%%%%%%%%%%%%%%%%%%%%%%%%%%%%%%%%%%%%%%%%%%%

\end{document}